\documentclass{article}
\usepackage[backend=biber]{biblatex}
\usepackage{graphicx}
\usepackage[colorlinks=true]{hyperref}
\usepackage{booktabs}
\usepackage{siunitx}
\usepackage[]{amsmath}
\usepackage{gensymb}
\usepackage{mathtools}
\usepackage{fancyref}
\usepackage{chemfig}
\addbibresource{/home/giorgio/Bibliography/bibliography.bib}
\hypersetup{
    colorlinks=true,
    linkcolor=blue,
    filecolor=magenta,  
    citecolor=blue,    
    urlcolor=cyan,
    pdftitle={Selective Laser Synthering for renewable polymers},
    bookmarks=true,
}

\author{Giorgio De Trane}
\title{\textbf{Selective Laser Synthering for renewable polymers}}

\begin{document}
    \setlength{\parindent}{0pt}
    \maketitle
    \begin{center}
        \includegraphics[width=0.8\textwidth]{Pictures/polito_logo.eps}\linebreak\newline
        \textit{Year 2021/2022}
    \end{center}

    \newpage
    \tableofcontents
    \newpage
    \section{Introduction\label{Intro}}

    \textit{Additive manufacturing} (AM) is a broad term that encompasses several manufacturing techniques, characterized by their additive nature, as the name suggests, 
    in contrast with more traditional subtractive processes.
    
    AM techniques are applied to a vast range of materials, including ceramics, polymers and metal alloys, some of which are specifically developed or optimized to 
    these kinds of applications. \\ 
    
    The main advantage of AM is the ability to produce complex shapes in a relatively short time. 
    These geometries are either too hard or even impossible to reproduce with subtractive manufacturing techniques, which often require multiple steps, using different 
    pieces of equipment, trained personnel, etc. \\

    Given the same material, the complex shapes allowed by AM can replace components made of multiple assembled parts with a unique solid piece of comparable or even better mechanical properties. 
    
    The inherent flexibility of AM often allows product designers to simplify or even entirely bypass the very strict CAD workflow (which is intrinsically tied to 
    traditional manufacturing processes) and make use of organic and/or generative modeling. 

    As a consequence of better design choices and minimal need of post-processing of AM objects, far less raw material is wasted, compared to subtractive manufacturing techniques, leading to 
    long term lowering of costs, faster design-to-market pipelines and, last but not least, lower emissions and environmental impact \autocite*{Recent_progress_polymers_AM}. 

    \subsection{Polymers in Additive Manufacturing\label{Polymers_in_AM}}

    Polymers and their composite materials have been used in all sorts of fields, ranging from arts and crafts all the way to advanced biomedical and aerospace applications, thanks 
    to their unique and varied extended range of properties. \\  
    
    The rapid advancement of AM, where polymers have been extensively used for prototyping, in the form of resins, filaments, powders and viscous inks, 
    has increased the demand for high-performance polymers, in order to take advantage of their quicker printing times (compared to metals) as well as their lower cost, while still
    maintaining good mechanical properties for an end product, rather than just a prototype \autocite*{Recent_progress_polymers_AM}. \\ 

    The urge for drastically reducing the environmental impact of human activities involves every production field, including AM, which can be inherently less impactful
    than traditional manufacturing processes, given the same material and final product to achieve. 

    A consequence of the concerns about climate change and its potentially catastrophic outcomes is the research in the field of eco-friendly materials, including polymers 
    that could be used in AM. 
    
    A great example is \textit{PLA} (PolyLactic Acid), a polymer widely used in 3D printing, whose monomer is obtained by fermenting starches, such as corn starch. \\ 

    Many new eco-friendly polymers have been and are currently being studied for AM techniques, but this case study will focus mostly on materials that can be potentially turned 
    into powders for \textit{PBF} (Powder Bed Fusion) techniques or filaments for \textbf{ME} (Material Extrusion). 
    
    \subsection{Common AM techniques for polymers\label{AM_techniques_summary}}
    
    Polymers can be processed with several AM techniques, including, but not limited to: 

    \begin{itemize}
        \item \textbf{VP} (Vat Photopolymerization), which make use of UV lights (or other radiations) to solidify photosensitive resins. This class of AM processes can produce 
                        parts with the highest resolution among all AM methods \autocite*{Recent_progress_polymers_AM};
        \item \textbf{MJ} (Material Jetting), which consist of a deposition of viscous fluids (either in droplets or in a continuous fashion), solidified by different agents 
                            (time interacting chemicals, heating, cooling, drying, photopolymerization, etc.). These processes include several patented methods, characterized by 
                            high speed printing \autocite*{Recent_progress_polymers_AM};
        \item \textbf{PBF} (Powder Bed Fusion), where the object is printed by locally fusing a powder bed (with a pulsing energy source -such as lasers- or with a local deposition of chemicals),
                            layed out in a layer-by-layer fashion \autocite*{Recent_progress_polymers_AM};
        \item \textbf{ME} (Material Extrusion), where each layer is printed by direct deposition of materials through a nozzle, that solidify as they cool down \autocite*{Recent_progress_polymers_AM}. 
    \end{itemize}

    \subsubsection{FDM (Fused Deposition Modeling) \label{FDM_general}}

    FDM is a ME technique and the most well known AM process, commonly named 3D printing in popular media. \\
    
    The process consists of a direct layer-by-layer deposition of a plastic filament, heated up to its melting point and extruded through a nozzle. 
    
    This method has gained a lot of popularity in the last few years, given its general ease of use, low cost of both materials and equipment and its growing 
    community of enthusiasts. 








    \clearpage

    \printbibliography
    
\end{document}